% Source: https://artofproblemsolving.com/community/c3h2918242_argentina_nmo_round_2_intercollegiate_geo__level_1_200616
% Stable identity: aops-topic-2918242-q08
$CB=CD=CE$. In triangle $BCE$, $CB=CE$ and $\angle CBE=70^\circ-10^\circ=60^\circ$, so it is equilateral and $\angle BCE=60^\circ$. Then $\angle ECD=110^\circ$, and the isosceles triangle $CED$ has base angles $35^\circ$, so $\angle ADC=35^\circ$.
\begin{asy}
usepackage("amsmath");
size(3.5cm);
defaultpen(fontsize(10pt));
pair C=(0,0),B=dir(190),E=dir(130),D=dir(20),A=(-0.8112,0.8112);draw(A--B--C--D--cycle);draw(B--E);draw(C--E,dashed);
label("$A$",A,NW);label("$B$",B,W);label("$C$",C,S);label("$D$",D,E);label("$E$",E,N);
\end{asy}

At $E$: $\angle BEC+\angle CED=60^\circ+35^\circ=95^\circ$, so $\angle AEB=85^\circ$ and $\angle BAD=180^\circ-10^\circ-85^\circ=85^\circ$. Answer: $\boxed{\angle BAD=85^\circ,\ \angle ADC=35^\circ}$ (the angle sum $85+70+170+35=360$ checks).
